\section{Purpose}

\textbf{[STIPULATION]} This document addresses two related points: (1)
$E=mc^2$ is in natural units identical to $E=m$ -- the only difference lies in
the convention of setting $c$ as a constant. (2) The three-clock experiment
(following Sci.\ Rep.\ 2024) shows that all physical quantities -- lengths, masses
-- can be traced back to a single time standard, which metrologically supports the
T0 duality $\tilde T\cdot m=1$.

\section{E = mc² = E = m}

\textbf{[B]} In natural units ($c=1$):
\[
  E = mc^2 = m\times1^2 = m.
\]
This is not an approximation -- it is the same equation. The difference between
Einstein and T0 lies not in the physics but in the convention.

\textbf{[B]} What is $c$ really? $c=L/T$ is a ratio of length to time. Einstein's
setting $c=299\,792\,458\,\text{m/s}$ freezes this ratio and makes the equation
$E=mc^2$ necessary. T0 theory allows $c$ to be a dynamic ratio -- then $E=m$ is
the more fundamental form. Both formulations are physically equivalent; what
differs is the ontological status of $c$.

\textbf{[B]} \textbf{The constant illusion.} Setting $c=\text{const.}$ is a
legitimate normalisation choice, not a constant of nature. With it:
\[
  E=mc^2\;\text{(in SI units)}
  \quad\equiv\quad
  E=m\;\text{(in natural units)}.
\]
The apparent complexity of $c^2$ is a consequence of the unit choice, not of the
physics. The same holds for $\hbar$ and $k_B$ (A085).

\textbf{[S]} \textbf{The time dilation paradox.} In the standard view, time is
relative, mass is constant. In T0, time is absolute, mass is variable. Both views
describe the same measurements; which to prefer remains an ontological choice (A060).

\section{The Three-Clock Experiment (Sci. Rep. 2024)}

\textbf{[B]} The experiment (Sci.\ Rep.\ 2024, DOI: 10.1038/s41598-024-71907-0)
shows: a single time standard suffices as a starting point to define and measure all
physical quantities -- time intervals, lengths, masses.

\textbf{[B]} \textbf{Length measurement from time.} Lengths are determined
exclusively from time differences: $L=c\cdot\Delta t$. This sets $c=1$ (convention)
and gives lengths in units of time -- exactly what the T0 duality requires.

\textbf{[B]} \textbf{Mass measurement from frequencies.} Via the Compton wavelength
$\lambda_C=\hbar/(mc)$ and metrological procedures (Kibble balance), masses can be
traced to the time standard: $m\propto1/T$ with $T$ as the Compton period. This is
the metrological realisation of $\tilde T\cdot m=1$.

\textbf{[K]} \textbf{Connection to T0 duality.} The reduction to a single time
standard is not new -- T0 already achieves this in the units documents (A080/A085).
The three-clock experiment is an external, independent confirmation of the metrological
consistency of this reduction. It shows that the programme is technically feasible,
but presupposes the same convention choice ($c=1$).

\section{Reference-Point Revolution}

\textbf{[S]} T0 theory makes the implicit reference-point shift explicit:
\begin{center}
\renewcommand{\arraystretch}{1.3}
{\small
\begin{tabular}{lll}
\toprule
Epoch & Reference & Consequence \\
\midrule
Antiquity & Earth as centre & Epicycles \\
Newton & Sun as reference & Kepler's laws \\
Einstein & $c=\text{const}$, space as reference & Time dilation \\
T0 & Time as reference & Mass variation \\
\bottomrule
\end{tabular}
}
\renewcommand{\arraystretch}{1.0}
\end{center}
Every shift is a convention choice -- none is ``more correct'' than another.
What counts: which describes the physics more simply.

\section{Verification Script}

\begin{itemize}[leftmargin=2em,itemsep=2pt,topsep=2pt]
  \item Numerical equivalence $E=mc^2=E=m$ for various $c$ values; length and
    mass measurement from a single time standard:
    \texttt{python/A\_Serie\_Skripte/a262\_c\_konvention.py}
\end{itemize}

\section{Open Questions}

\textbf{The ontological question is not decided.} Which view is preferable --
time relative and mass constant, or time absolute and mass variable -- is a
question of interpretation, not physics. The A-series makes no judgement (A060).

\textbf{The three-clock experiment does not cover the full T0 claim.} It confirms
the metrological reduction; it does not confirm the geometric derivation of $\xi$
or the mass formulas.

\section{Supersedes}

This document subsumes: Dok.~077 ($E=mc^2=E=m$, $c$-convention, time dilation
paradox, reference-point revolution), Dok.~029 (three-clock experiment Sci.\ Rep.\
2024, length and mass measurement from time standard, connection to T0 duality).

\section{Use of AI Tools}

This document was created with the assistance of AI tools. Responsibility
for content and errors rests with the author. Full wording of the rules in A010.
